\section{Signed finite-radius envelopes}\label{sec:assembly}

The common prices control unfavorable directional components. The
direction law controls the endpoints created when a direction changes.
We now combine those estimates, keeping the remaining directional
projections signed until integration by parts. The essential cancellation
is that the derivatives of a probability row sum to zero. Subtracting
the potential of the preceding history then leaves only oscillations
at smaller radii.

Throughout this section, $g\in C_c^\infty(\R^2)$ is real,
$u_t=A_tg$, $V=\int|\nabla g|$, and $r_k=2^{-k}$ for
$1\le k\le m$. Use the prices and label law of
Sections~\ref{sec:prices} and~\ref{sec:labels}. For $1\le i\le j\le m$,
put
\begin{equation}\label{eq:run-envelope}
 F_{ij}(x)=\max_{r_j\le t\le2r_i}u_t(x).
\end{equation}
The maximum is attained, since $t\mapsto u_t(x)$ is continuous.
These envelopes are compactly supported: all their radii are at most
$2r_1=1$, so all their averages vanish outside the closed
$1$-neighborhood of $\supp g$. They are Lipschitz, with common
Lipschitz bound $\|\nabla g\|_\infty$.

Outside one null set all the finitely many $F_{ij}$ are
differentiable. At any such point $x$, every radius $t$ attaining
$F_{ij}(x)$ satisfies
\begin{equation}\label{eq:envelope-touching}
 \nabla F_{ij}(x)=\nabla u_t(x).
\end{equation}
Indeed, with this radius held fixed, the differentiable function
$F_{ij}-u_t$ is nonnegative everywhere and is zero at $x$. Its
gradient there is zero. This proves the assertion also when there
are several attaining radii. The same touching principle underlies
the optimizing-radius identity of Haj{\l}asz and Mal\'y
\cite[Theorem~2 and Remark~3]{HajlaszMaly2010}; here it applies
directly to the smooth signed averages. No selected radius is
differentiated.

To prove Proposition~\ref{prop:finite-band}, it suffices first to
bound the normalized band by an absolute constant independent of $m$:
\begin{equation}\label{eq:finite-band}
 \int_{\R^2}|\nabla F_{1m}|\,\dd x\le C\int_{\R^2}|\nabla g|\,\dd x.
\end{equation}

\begin{proof}[Proof of Proposition~\ref{prop:finite-band}]
We first bound the normalized band in \eqref{eq:finite-band}.
For a label history $\mathbf n=(n_1,\ldots,n_m)\in\mathcal N^m$,
a \emph{run} is a maximal interval $[i,j]$ of indices on which
the label is constant. Its radius interval is $[r_j,2r_i]$.
Adjacent intervals meet at an endpoint, and their union is
$[r_m,1]$. Thus
\[
 F_{1m}=\max_{[i,j]\text{ a run}}F_{ij}.
\]
At a point where all the envelopes are differentiable, some run
attains this maximum, and its gradient agrees with $\nabla F_{1m}$
by the same touching argument.

The cone inequality \eqref{eq:cone-algebra} states that, for every
$n\in\mathcal N$ and $z\in\R^2$,
\[
 |z|\le C\left[n\cdot z+
                 \sum_{\pm}(-e_\pm(n)\cdot z)_+\right],
\]
and the bracket is nonnegative. Apply it to a run attaining
$F_{1m}$, add the nonnegative brackets from the other runs, and
average over the history law at $x$. This gives
\begin{equation}\label{eq:run-master}
 |\nabla F_{1m}|
 \le C\,\E_x\sum_{[i,j]\text{ a run}}
 \left[n_i\cdot\nabla F_{ij}
       +\sum_{\pm}(-e_\pm(n_i)\cdot\nabla F_{ij})_+\right]
 \quad\text{a.e.}
\end{equation}
There are finitely many histories and runs. In particular, their
expectations and all differentiations of their weights below are
finite sums.

\paragraph{The unfavorable components and the ends of runs.}
For a run $[i,j]$, choose any attaining radius $t$ at the point
under consideration, and a bin $[r_k,2r_k]$ containing it, with
$i\le k\le j$. If $i<k<j$, the interval
$[t-r_k/4,t+r_k/4]$ lies inside $[r_j,2r_i]$. The maximizing
property therefore gives $t\in\mathcal I_k(x)$. If $k=i$ or
$k=j$, the full bin is still contained in the run interval, so
$t$ maximizes on that bin and belongs to $\mathcal J_k(x)$.
Touching, including at tied radii, identifies the gradients.

Write
\[
 B_k(x)=\sup_{t\in\mathcal J_k(x)}|\nabla u_t(x)|,
 \qquad
 A_{k,n}(x)=\sup_{t\in\mathcal I_k(x)}
                  \sum_{\pm}(-e_\pm(n)\cdot\nabla u_t(x))_+.
\]
A bin is an end bin of its run only if a run starts at $k$ or
ends at $k$. These events have probabilities $R_k$ and $R_{k+1}$,
respectively, with the sentinel conventions $R_1=R_{m+1}=1$.
Also each bin belongs to exactly one run. Since each of the two
negative projections is bounded by the gradient norm, we obtain
\begin{equation}\label{eq:negative-sum}
 \E_x\sum_{[i,j]\text{ a run}}\sum_\pm
       (-e_\pm(n_i)\cdot\nabla F_{ij})_+
 \le\sum_{k=1}^m\sum_{n\in\mathcal N}
          \Prob_x(n_k=n)A_{k,n}(x)
       +2\sum_{k=1}^m(R_k+R_{k+1})B_k(x).
\end{equation}
No measurable choice of the attaining radius is needed: this is
a pointwise bound by the measurable suprema on the right.

Here are the spatial sums in this estimate. For a level-$k$ cell
$Q=Q_p$ and $q\in\mathcal R(Q)$, every $x\in Q$ satisfies
$|x-q|_\infty\le21r_k$. Thus $Q\subset D(q)$, and $Q$ also
lies in the $50r_k$ neighborhood used in Lemma~\ref{lem:good-label}.
If $(q,n)$ is bad and $E_k(q)\le c_*$, that lemma says
$\Prob_x(n_k=n)=0$ throughout $Q$. Every remaining bad-root
charge is bounded by
$a_q\le c_*^{-1}a_qE_k(q)$. Take the supremum of each label
probability on $Q$ and apply the directional part of
Proposition~\ref{prop:prices}. Its price charges sum by
\eqref{eq:price-budget}; its root charges sum by
\eqref{eq:neighbor-budget}. The finite palette and the $41^2$
cells having a given vertex in their root neighborhood contribute
only absolute factors.

For the last sum in \eqref{eq:negative-sum}, the endpoint estimate
in Proposition~\ref{prop:prices} gives, cell by cell,
\[
 \int_Q(R_k+R_{k+1})B_k
 \le C\sup_Q(R_k+R_{k+1})
                \left(D_Q+\sum_{q\in\mathcal R(Q)}a_q\right).
\]
The price part sums because $R_k+R_{k+1}\le2$.
For each root part, enlarge $Q$ to $D(q)$ and use the switch
part of \eqref{eq:label-packing}, again with bounded cell
multiplicity. We have proved
\begin{equation}\label{eq:negative-integral}
 \int_{\R^2}\E_x\sum_{[i,j]\text{ a run}}\sum_\pm
       (-e_\pm(n_i)\cdot\nabla F_{ij})_+\,\dd x\le CV.
\end{equation}

\paragraph{Baselines and weighted integration by parts.}
It remains to integrate the signed projections in
\eqref{eq:run-master}. Subtract $u_{r_i}$ from the envelope of
every run starting at $i$. Since $r_i$ is a forced endpoint in
$\mathcal J_i(x)$, the same endpoint and packing estimates give
\begin{equation}\label{eq:baseline-cost}
 \int\E_x\sum_{[i,j]\text{ a run}}|\nabla u_{r_i}|
 =\sum_{i=1}^m\int R_i|\nabla u_{r_i}|
 \le\sum_{i=1}^m\int(R_i+R_{i+1})B_i\le CV.
\end{equation}
For a fixed history define its vector potential by
\begin{equation}\label{eq:path-potential}
 \Phi_{\mathbf n}(x)=
       \sum_{[i,j]\text{ a run}}n_i\bigl(F_{ij}(x)-u_{r_i}(x)\bigr).
\end{equation}
Its run partition and labels are fixed; only its scalar functions
depend on $x$. It is a compactly supported Lipschitz vector field,
and its divergence is the sum of the signed projections after
baseline subtraction.

Use the auxiliary state $n_0=*$ to write the initial row as
$\pi_1(x;*,n_1)$. The probability of the history is
\[
 w_{\mathbf n}(x)=\prod_{s=1}^m\pi_s(x;n_{s-1},n_s).
\]
Every factor is Lipschitz and takes values in $[0,1]$. The weak
product rule and compact support of $w_{\mathbf n}\Phi_{\mathbf n}$
therefore give
\begin{equation}\label{eq:weighted-ibp}
 \sum_{\mathbf n}\int w_{\mathbf n}\,
                       \operatorname{div}\Phi_{\mathbf n}\,\dd x
 =-\sum_{\mathbf n}\int
                       \Phi_{\mathbf n}\cdot\nabla w_{\mathbf n}\,\dd x.
\end{equation}
Indeed the integral of the divergence of a compactly supported
Lipschitz vector field is zero, as follows by testing against a
smooth function equal to one on its support. We must exploit the
sign of the right side before estimating its size.

\paragraph{Cancellation after a fixed prefix.}
Differentiate the product defining $w_{\mathbf n}$ directly. For
its term at step $k$, fix the prefix
$\mathbf a=(n_1,\ldots,n_{k-1})$ and set
\[
 W_{\mathbf a}=\prod_{s<k}\pi_s(x;n_{s-1},n_s),
 \qquad
 J_{k,\mathbf n}=\prod_{s>k}\pi_s(x;n_{s-1},n_s).
\]
Empty products equal one. The corresponding term of
$\nabla w_{\mathbf n}$ is
$W_{\mathbf a}\nabla\pi_k(x;n_{k-1},n_k)J_{k,\mathbf n}$.
This formula remains valid when some transition entries vanish;
it involves no division by them.

For each fixed prefix and $n_k$, summing the future product gives
\begin{equation}\label{eq:future-normalization}
 \sum_{n_{k+1},\ldots,n_m}J_{k,\mathbf n}=1.
\end{equation}
To see this, first sum in $n_m$ using the normalization of the last
row, and repeat backward to $n_{k+1}$. Also, wherever the rows are
differentiable,
\begin{equation}\label{eq:row-derivative-zero}
 \sum_{n_k\in\mathcal N}\nabla\pi_k(x;n_{k-1},n_k)=0,
\end{equation}
by differentiating their finite sum, which is the constant one.

Let $\Psi_{\mathbf a}$ be the potential of this prefix with its
last run closed at $k-1$: explicitly, use the sum
\eqref{eq:path-potential} for the runs of $(n_1,\ldots,n_{k-1})$.
For $k=1$ set $\Psi_{\varnothing}=0$. This potential is
independent of $n_k$ and of every future label. Equations
\eqref{eq:future-normalization}--\eqref{eq:row-derivative-zero}
therefore imply the exact identity
\begin{align}
 &\sum_{n_k,\ldots,n_m}
        J_{k,\mathbf n}\Phi_{\mathbf n}\cdot
                  \nabla\pi_k(x;n_{k-1},n_k)\notag\\
 &\hspace{8mm}=
   \sum_{n_k,\ldots,n_m}J_{k,\mathbf n}
        (\Phi_{\mathbf n}-\Psi_{\mathbf a})\cdot
                  \nabla\pi_k(x;n_{k-1},n_k).
 \label{eq:prefix-subtraction}
\end{align}
The subtraction precedes every absolute value in the next estimate.

We claim that each continuation of the prefix satisfies
\begin{equation}\label{eq:prefix-potential}
 |\Phi_{\mathbf n}(x)-\Psi_{\mathbf a}(x)|
       \le\sum_{l=k}^m v_l(x),
 \qquad v_l=\sup_{0<t,t'\le4r_l}|u_t-u_{t'}|.
\end{equation}
All already completed runs cancel. If the last prefix run
$[i,k-1]$ continues to $j\ge k$, its baseline cancels as well,
leaving $n_i(F_{ij}-F_{i,k-1})$. The old radius interval contains
$r_{k-1}$, and the extension adds only radii below that number.
Consequently
\[
 0\le F_{ij}-F_{i,k-1}
 \le\sup_{0<t\le r_{k-1}}(u_t-u_{r_{k-1}})_+
 \le v_k,
 \qquad r_{k-1}=2r_k.
\]
Every new run $[l,j]$, beginning at $l\ge k$, contributes a vector
of norm $F_{lj}-u_{r_l}\le v_l$, because its envelope includes
the baseline and all its radii are at most $2r_l$. If the prefix
run continues, the new runs start strictly after $k$; if there
is a switch at $k$, the closed prefix itself is unchanged. The
empty-prefix case consists entirely of new runs. Since their
starting indices are distinct, these observations prove
\eqref{eq:prefix-potential}. In particular, the size of an old
large-radius envelope never enters the bound.

Now take absolute values in \eqref{eq:prefix-subtraction} and sum
over prefixes. The future products are nonnegative and sum to
one. The sum of the norms of the step-$k$ product-rule terms is
exactly
\begin{align*}
 &\sum_{\mathbf a}W_{\mathbf a}
     \sum_{n_k}|\nabla\pi_k(x;n_{k-1},n_k)|
       \sum_{n_{k+1},\ldots,n_m}J_{k,\mathbf n}\\
 &\qquad=
   \sum_{n\in\mathcal N}\zeta_{k-1}(x;n)
                  \sum_{n'}|\nabla\pi_k(x;n,n')|
   =\frac{T_k(x)}{r_k}\quad(k\ge2).
\end{align*}
For $k=1$ the same identity uses the single initial row, exactly
as in \eqref{eq:labels-derivative-costs}. Combining these identities
with the product expansion and \eqref{eq:weighted-ibp} yields
\begin{equation}\label{eq:flux-cost}
 \left|\sum_{\mathbf n}\int
                w_{\mathbf n}\operatorname{div}\Phi_{\mathbf n}\right|
 \le\sum_{1\le k\le l\le m}
                    \int\frac{T_k(x)}{r_k}v_l(x)\,\dd x.
\end{equation}

\paragraph{Paying for the smaller-radius oscillations.}
On a level-$l$ cell, the value estimate in
Proposition~\ref{prop:prices} gives
\[
 \int_Qv_l\le Cr_l
             \left(D_Q+\sum_{q\in\mathcal R(Q)}a_q\right).
\]
Multiply this by $r_k^{-1}\operatorname*{ess\,sup}_QT_k$.
For the price terms, $T_k\le C$ and
\[
 \sum_{k=1}^l\frac{r_l}{r_k}
       =\sum_{k=1}^l2^{k-l}\le2.
\]
Their full sum is therefore at most $C\sum_{l,Q}D_Q\le CV$.
For the root terms use $Q\subset D(q)$ and reorder the finite
nonnegative sums. They are at most an absolute multiple of
\[
 \sum_{l=1}^m\sum_{q\in V_l}a_q
       \sum_{k=1}^l2^{k-l}
                   \operatorname*{ess\,sup}_{D(q)}T_k\le CV
\]
by \eqref{eq:label-packing}. Thus \eqref{eq:flux-cost} is at
most $CV$. Together with \eqref{eq:baseline-cost}, this bounds
the integral of the signed projections in \eqref{eq:run-master}
from above by $CV$. Equation~\eqref{eq:negative-integral} pays
for the remaining terms and proves \eqref{eq:finite-band}.

For $b/a=2^m$ with $m\ge1$, apply the normalized result to
$g_b(x)=g(bx)$. Its normalized envelope is
$S_{a,b}g(bx)$, and both gradient integrals acquire the same
factor $b^{-1}$ under the change of variables. This proves the
stated bound for all bands in the proposition. Finally, when $a=b$, convolution
gives $\nabla A_ag=A_a(\nabla g)$, and Tonelli's theorem yields
$\int|\nabla A_ag|\le\int|\nabla g|$.
\end{proof}
